% ==========================================================================================
% Chapter 4: Website Usage
% ==========================================================================================

บทนี้อธิบายวิธีการใช้งานระบบ GymBro Management System อย่างละเอียด แบ่งตามสิทธิ์การใช้งานของผู้ใช้ 3 กลุ่มหลัก ได้แก่ ผู้ใช้งานทั่วไป (Public), สมาชิก (Member), และผู้ดูแลระบบ (Admin) โดยมีการแสดงภาพประกอบหน้าจอทั้งในโหมดปกติ (Light Mode) และโหมดมืด (Dark Mode)

\section{ส่วนผู้ใช้งานทั่วไป (Public Section)}
ส่วนนี้สำหรับผู้ที่ยังไม่ได้เข้าสู่ระบบ สามารถใช้งานได้ทุกคน

\subsection{การใช้งานหน้าเข้าสู่ระบบ (Login Page)}
หน้าแรกของระบบที่ใช้สำหรับยืนยันตัวตนก่อนเข้าใช้งาน

\subsubsection{ขั้นตอนการเข้าสู่ระบบ}
\begin{enumerate}
    \item เปิดเว็บเบราว์เซอร์และเข้าสู่หน้าเว็บไซต์ ระบบจะแสดงแบบฟอร์ม Login
    \item กรอก \textbf{Email Address} ที่ได้ลงทะเบียนไว้
    \item กรอก \textbf{Phone Number} (ใช้แทนรหัสผ่าน)
    \item กดปุ่ม \textbf{SIGN IN}
    \item หากข้อมูลถูกต้อง ระบบจะตรวจสอบสิทธิ์ (Role) และเปลี่ยนหน้าไปยัง Dashboard ที่เหมาะสม (User หรือ Admin)
    \item หากข้อมูลไม่ถูกต้อง ระบบจะแสดงข้อความแจ้งเตือนสีแดง "Invalid credentials"
\end{enumerate}

\begin{figure}[htbp]
    \centering
    \includegraphics[width=0.9\textwidth]{Image/login.png}
    \caption{หน้าจอเข้าสู่ระบบ (Login) ประกอบด้วยช่องกรอกอีเมลและเบอร์โทรศัพท์ พร้อมปุ่มเข้าสู่ระบบ}
    \label{fig:login}
\end{figure}

\subsection{การใช้งานหน้าลงทะเบียน (Register Page)}
สำหรับผู้ใช้งานใหม่ที่ต้องการสมัครสมาชิก

\subsubsection{ขั้นตอนการสมัครสมาชิก}
\begin{enumerate}
    \item ในหน้า Login ให้กดลิงก์ "Don't have an account? Register here"
    \item ระบบจะพาไปที่หน้า Register
    \item กรอก \textbf{Full Name} (ชื่อ-นามสกุล)
    \item กรอก \textbf{Email Address} (ห้ามซ้ำกับที่มีในระบบ)
    \item กรอก \textbf{Phone Number} (ใช้สำหรับล็อกอินในภายหลัง)
    \item กดปุ่ม \textbf{CREATE ACCOUNT}
\end{enumerate}

\begin{figure}[htbp]
    \centering
    \includegraphics[width=0.9\textwidth]{Image/register.png}
    \caption{หน้าจอลงทะเบียน (Register) สำหรับกรอกข้อมูลส่วนตัวเพื่อสมัครสมาชิกใหม่}
    \label{fig:register}
\end{figure}

\subsubsection{การรอการอนุมัติ (Pending Approval)}
เมื่อสมัครสมาชิกสำเร็จ ระบบจะแสดงข้อความแจ้งเตือนว่าบัญชีอยู่ในสถานะ \textbf{รอการอนุมัติ (Pending)} ผู้ใช้งานจะต้องติดต่อเจ้าหน้าที่เพื่อชำระค่าบริการ และรอให้เจ้าหน้าที่ยืนยันสถานะบัญชีจึงจะสามารถเข้าใช้งานได้

\begin{figure}[htbp]
    \centering
    \includegraphics[width=0.9\textwidth]{Image/registered_alert.png}
    \caption{หน้าต่างแจ้งเตือนการสมัครสมาชิกสำเร็จและรอการอนุมัติจากผู้ดูแลระบบ}
    \label{fig:registered_alert}
\end{figure}

\newpage
\section{ส่วนสมาชิก (Member Section)}
ส่วนนี้สำหรับผู้ใช้งานที่ล็อกอินด้วยสิทธิ์สมาชิก (User)

\subsection{การใช้งานหน้าแดชบอร์ดสมาชิก (User Dashboard)}
หน้าหลักของสมาชิก ใช้สำหรับดูตารางฝึกซ้อมและจองเซสชัน

\subsubsection{การดูตารางฝึกซ้อม}
\begin{enumerate}
    \item เมื่อเข้าสู่ระบบสำเร็จ จะพบตารางแสดงรายการจองทั้งหมดในวันนี้ (All Bookings)
    \item รายการที่เป็นของตนเอง (My Session) จะถูกไฮไลต์ด้วยสีพื้นหลังที่แตกต่างเพื่อให้สังเกตได้ง่าย
    \item ตารางแสดงข้อมูลเวลา (Time), กิจกรรม/อุปกรณ์ (Activity), และชื่อเทรนเนอร์ (Trainer)
\end{enumerate}

\begin{figure}[htbp]
    \centering
    \includegraphics[width=0.9\textwidth]{Image/user_customer_schedule.png}
    \caption{หน้าแดชบอร์ดสมาชิกแสดงตารางการฝึกซ้อม (User Schedule) พร้อมไฮไลต์รายการจองของตนเอง}
    \label{fig:user_dashboard_schedule}
\end{figure}

\subsection{การจองเซสชัน (Booking Session)}

\subsubsection{ขั้นตอนการจอง}
\begin{enumerate}
    \item กดปุ่ม \textbf{BOOK SESSION} ที่มุมขวาบน
    \item ระบบจะแสดงหน้าต่าง (Modal) สำหรับกรอกข้อมูลการจอง
    \item เลือก \textbf{Date} (วันที่), \textbf{Time} (เวลา), และ \textbf{Duration} (ระยะเวลา)
    \item ระบบจะตรวจสอบความว่าง (Availability) ของเทรนเนอร์และอุปกรณ์โดยอัตโนมัติ หากไม่ว่างตัวเลือกนั้นจะถูกปิดการใช้งาน (Disabled)
    \item เลือก \textbf{Trainer} และ \textbf{Equipment} ที่ต้องการ
    \item กดปุ่ม \textbf{CONFIRM BOOKING}
\end{enumerate}

\begin{figure}[htbp]
    \centering
    \includegraphics[width=0.9\textwidth]{Image/user_customer_booking.png}
    \caption{หน้าต่างสำหรับการจองเซสชัน (Booking Modal) ให้ผู้ใช้เลือกวัน เวลา เทรนเนอร์ และอุปกรณ์}
    \label{fig:user_booking_modal}
\end{figure}

\subsubsection{การยืนยันการจอง}
เมื่อทำการจองสำเร็จ ระบบจะแสดงข้อความแจ้งเตือน (Alert) เพื่อยืนยันว่าการจองได้รับการบันทึกเรียบร้อยแล้ว และตารางจะถูกรีเฟรชเพื่อแสดงรายการจองใหม่

\begin{figure}[htbp]
    \centering
    \includegraphics[width=0.9\textwidth]{Image/customer_booked.png}
    \caption{หน้าจอแสดงสถานะการจองสำเร็จ (Booking Success) พร้อมแสดงรายการใหม่ในตาราง}
    \label{fig:customer_booked}
\end{figure}

\begin{figure}[htbp]
    \centering
    \includegraphics[width=0.9\textwidth]{Image/customer_booked_alert.png}
    \caption{ข้อความแจ้งเตือนยืนยันการจองสำเร็จ (Booking Alert) ที่มุมขวาบนของหน้าจอ}
    \label{fig:customer_booked_alert}
\end{figure}

\subsection{การใช้งานหน้าข้อมูลส่วนตัว (My Profile)}
หน้าแสดงรายละเอียดข้อมูลสมาชิก และสถานะบัญชี (แสดงตัวอย่างในโหมดมืด)

\begin{enumerate}
    \item ที่แถบเมนูด้านซ้าย เลือกเมนู \textbf{My Profile}
    \item ระบบจะแสดงข้อมูลส่วนตัว ได้แก่ ชื่อ-นามสกุล, อีเมล, และเบอร์โทรศัพท์
    \item แสดงประเภทสมาชิก (Member Type) และวันหมดอายุสมาชิก
    \item แสดงสถานะบัญชี (Account Status)
\end{enumerate}

\begin{figure}[htbp]
    \centering
    \includegraphics[width=0.9\textwidth]{Image/dark_user_profile.png}
    \caption{หน้าข้อมูลส่วนตัว (My Profile) ในโหมดมืด แสดงรายละเอียดสมาชิกและสถานะบัญชี}
    \label{fig:dark_user_profile}
\end{figure}

\newpage
\section{ส่วนผู้ดูแลระบบ (Admin Section)}
ส่วนนี้สำหรับผู้ใช้งานที่ล็อกอินด้วยสิทธิ์ผู้ดูแลระบบ (Admin)

\subsection{การเข้าสู่ระบบสำหรับผู้ดูแลระบบ (Admin Login)}
ผู้ดูแลระบบสามารถเข้าสู่ระบบผ่านหน้า Login เดียวกับผู้ใช้งานทั่วไป หรือผ่าน URL เฉพาะ

\begin{figure}[htbp]
    \centering
    \includegraphics[width=0.9\textwidth]{Image/admin_login.png}
    \caption{หน้าจอเข้าสู่ระบบสำหรับผู้ดูแลระบบ (Admin Login) ใช้แบบฟอร์มเดียวกับผู้ใช้งานทั่วไป}
    \label{fig:admin_login}
\end{figure}

\subsection{การใช้งานหน้าแดชบอร์ดผู้ดูแลระบบ (Admin Dashboard)}
หน้าสรุปภาพรวมของระบบ แสดงสถิติและตารางงานวันนี้

\begin{figure}[htbp]
    \centering
    \includegraphics[width=0.9\textwidth]{Image/admin_dashboard.png}
    \caption{หน้าแดชบอร์ดผู้ดูแลระบบ (Admin Dashboard) แสดงจำนวนลูกค้า เทรนเนอร์ อุปกรณ์ และเซสชันวันนี้}
    \label{fig:admin_dashboard}
\end{figure}

\subsection{การจัดการข้อมูลลูกค้า (Customer Management)}
เมนู \textbf{Customers} สำหรับบริหารจัดการข้อมูลสมาชิก ตรวจสอบสถานะ และอนุมัติสมาชิกใหม่

\begin{figure}[htbp]
    \centering
    \includegraphics[width=0.9\textwidth]{Image/admin_member_management.png}
    \caption{หน้าจัดการข้อมูลลูกค้า (Customer Management) แสดงรายชื่อสมาชิกและปุ่มสำหรับจัดการข้อมูล}
    \label{fig:admin_customers}
\end{figure}

\subsection{การจัดการข้อมูลเทรนเนอร์ (Trainer Management)}
เมนู \textbf{Trainers} สำหรับบริหารจัดการข้อมูลผู้ฝึกสอน ระบุความเชี่ยวชาญ

\begin{figure}[htbp]
    \centering
    \includegraphics[width=0.9\textwidth]{Image/admin_trainer_management.png}
    \caption{หน้าจัดการข้อมูลเทรนเนอร์ (Trainer Management) แสดงรายชื่อและความเชี่ยวชาญของเทรนเนอร์}
    \label{fig:admin_trainers}
\end{figure}

\subsection{การจัดการอุปกรณ์ (Equipment Management)}
เมนู \textbf{Equipment} สำหรับบริหารจัดการอุปกรณ์และกำหนดสถานะการใช้งาน

\begin{figure}[htbp]
    \centering
    \includegraphics[width=0.9\textwidth]{Image/admin_equipment_management.png}
    \caption{หน้าจัดการอุปกรณ์ (Equipment Management) แสดงรายการอุปกรณ์และสถานะความพร้อมใช้งาน}
    \label{fig:admin_equipments}
\end{figure}

\subsection{การจัดการตารางฝึกซ้อม (Session Management)}
เมนู \textbf{Sessions} สำหรับดูและจัดการการจองทั้งหมด

\begin{figure}[htbp]
    \centering
    \includegraphics[width=0.9\textwidth]{Image/admin_training_sessions.png}
    \caption{หน้าจัดการตารางฝึกซ้อม (Session Management) แสดงรายการจองทั้งหมดและฟังก์ชันการค้นหา}
    \label{fig:admin_sessions}
\end{figure}

\subsection{การดูบันทึกระบบ (System Logs)}
เมนู \textbf{System Logs} สำหรับตรวจสอบประวัติการใช้งาน

\begin{figure}[htbp]
    \centering
    \includegraphics[width=0.9\textwidth]{Image/admin_report_syslog.png}
    \caption{หน้าบันทึกระบบ (System Logs) แสดงประวัติการกระทำต่างๆ ของผู้ใช้ในระบบ}
    \label{fig:admin_logs}
\end{figure}

\subsection{การออกรายงาน (Reports)}
ระบบมีฟังก์ชันสำหรับออกรายงานสรุปข้อมูลต่างๆ

\subsubsection{รายงานข้อมูลลูกค้าและเทรนเนอร์}

\begin{figure}[htbp]
    \centering
    \includegraphics[width=0.9\textwidth]{Image/admin_report_customer.png}
    \caption{รายงานข้อมูลลูกค้า (Customer Report) แสดงในรูปแบบตารางพร้อมสำหรับพิมพ์}
    \label{fig:admin_report_customer}
\end{figure}

\begin{figure}[htbp]
    \centering
    \includegraphics[width=0.9\textwidth]{Image/admin_report_trainer.png}
    \caption{รายงานข้อมูลเทรนเนอร์ (Trainer Report) แสดงรายชื่อและความเชี่ยวชาญ}
    \label{fig:admin_report_trainer}
\end{figure}

\subsubsection{การพิมพ์รายงาน (Printing Report)}

\begin{figure}[htbp]
    \centering
    \includegraphics[width=0.9\textwidth]{Image/admin_printing_report.png}
    \caption{หน้าตัวอย่างก่อนพิมพ์รายงาน (Print Preview) จัดรูปแบบหน้ากระดาษอัตโนมัติ}
    \label{fig:admin_printing_report}
\end{figure}

\newpage
\section{การแสดงผลในโหมดมืด (Dark Mode Interface)}
ระบบรองรับการแสดงผลในโหมดมืด (Dark Mode) เพื่อถนอมสายตาและประหยัดพลังงาน โดยมีการปรับเปลี่ยนโทนสีของพื้นหลัง ตัวอักษร และองค์ประกอบต่างๆ ให้เหมาะสมกับการใช้งานในที่ที่มีแสงน้อย

\subsection{หน้าจอผู้ดูแลระบบในโหมดมืด}

\begin{figure}[htbp]
    \centering
    \includegraphics[width=0.9\textwidth]{Image/dark_admin_dashboard.png}
    \caption{หน้าแดชบอร์ดผู้ดูแลระบบในโหมดมืด (Dark Admin Dashboard) ใช้โทนสีน้ำเงินเข้มและตัวอักษรสีอ่อน}
    \label{fig:dark_admin_dashboard}
\end{figure}

\begin{figure}[htbp]
    \centering
    \includegraphics[width=0.9\textwidth]{Image/dark_admin_member_management.png}
    \caption{หน้าจัดการข้อมูลลูกค้าในโหมดมืด (Dark Customer Management)}
    \label{fig:dark_admin_member_management}
\end{figure}

\begin{figure}[htbp]
    \centering
    \includegraphics[width=0.9\textwidth]{Image/dark_admin_trainer_management.png}
    \caption{หน้าจัดการข้อมูลเทรนเนอร์ในโหมดมืด (Dark Trainer Management)}
    \label{fig:dark_admin_trainer_management}
\end{figure}

\begin{figure}[htbp]
    \centering
    \includegraphics[width=0.9\textwidth]{Image/dark_admin_equipment_management.png}
    \caption{หน้าจัดการอุปกรณ์ในโหมดมืด (Dark Equipment Management)}
    \label{fig:dark_admin_equipment_management}
\end{figure}

\begin{figure}[htbp]
    \centering
    \includegraphics[width=0.9\textwidth]{Image/dark_admin_training_session.png}
    \caption{หน้าจัดการตารางฝึกซ้อมในโหมดมืด (Dark Session Management)}
    \label{fig:dark_admin_training_session}
\end{figure}

\begin{figure}[htbp]
    \centering
    \includegraphics[width=0.9\textwidth]{Image/dark_admin_report_syslog.png}
    \caption{หน้าบันทึกระบบในโหมดมืด (Dark System Logs)}
    \label{fig:dark_admin_report_syslog}
\end{figure}

\subsection{หน้ารายงานในโหมดมืด}

\begin{figure}[htbp]
    \centering
    \includegraphics[width=0.9\textwidth]{Image/dark_admin_report_customer.png}
    \caption{รายงานข้อมูลลูกค้าในโหมดมืด (Dark Customer Report)}
    \label{fig:dark_admin_report_customer}
\end{figure}

\begin{figure}[htbp]
    \centering
    \includegraphics[width=0.9\textwidth]{Image/dark_admin_report_trainer.png}
    \caption{รายงานข้อมูลเทรนเนอร์ในโหมดมืด (Dark Trainer Report)}
    \label{fig:dark_admin_report_trainer}
\end{figure}

\subsection{หน้าจอสมาชิกในโหมดมืด}

\begin{figure}[htbp]
    \centering
    \includegraphics[width=0.9\textwidth]{Image/dark_user_schedule.png}
    \caption{หน้าตารางฝึกซ้อมของสมาชิกในโหมดมืด (Dark User Schedule)}
    \label{fig:dark_user_schedule}
\end{figure}
